% Context from chapters/deep-learning.tex; for mathematical reference, not compilation.
\section{Deep Architectures}
\label{sec-deepnet-discr}

                                          
\subsection{Deep Network Structure}
\label{sec-deep-structure}

Deep learning constructs estimators $f(x,\be)$ by composing parameterized maps.
 
The hidden layers compute a feature representation of the input, and the output layer applies a linear or generalized linear predictor to those features. Unlike a fixed kernel feature map, this representation is learned from the training data.

   
\paragraph{Feedforward architectures.}

A feedforward chain has the structure introduced in~\eqref{eq-simple-lin-dag}, with the loss $\Ll$ omitted. Starting from $x_0=x$, compute
\eq{
	x_{\ell+1} = f_\ell(x_\ell,\be_\ell),  
}
so that $x_L=f(x,\be)$ with 
\eq{
	f(\cdot,\be) = f_{L-1}(\cdot,\be_{L-1}) \circ f_{L-2}(\cdot,\be_{L-2}) \circ \ldots \circ f_{0}(\cdot,\be_0)
}
where $\be=(\be_0,\ldots,\be_{L-1})$ collects the parameters, and each layer has the dimensions
\eq{
	f_{\ell}(\cdot,\be_\ell) : \RR^{n_\ell} \rightarrow \RR^{n_{\ell+1}}.
}
 
We focus on feedforward chains. More general architectures use branched computational graphs or share parameters across recurrent steps.

The parameters $\be$ are fitted by empirical risk minimization~\eqref{eq-erm-param}, typically using the stochastic methods of Section~\ref{sec-stochastic-optim}. The resulting objective is generally nonconvex, so Theorem~\ref{thm-conv-sgd} does not apply. Under suitable smoothness, noise, and step-size assumptions, one can instead bound measures of stationarity. Initialization and optimization dynamics can also provide implicit regularization.

For chain architectures, reverse mode computes the gradient of the ERM loss~\eqref{eq-grad-formula} as described in Section~\ref{sec-reversemode-simple}. For convolutional layers, it gives the backpropagation equations~\eqref{eq-backprop} below. More general computational graphs require the reverse pass to accumulate contributions from every branch and every use of a shared parameter; Section~\ref{sec-reverse-mode} explains this general procedure.

   
\paragraph{Deep MLP.}

A common computational block $f_\ell(\cdot,\be_\ell)$ consists of
\begin{rs}
	\item an affine map, $B_\ell \cdot + b_\ell$ with a matrix $B_\ell \in \RR^{\tilde n_{\ell} \times n_\ell}$ and a vector $b_\ell \in \RR^{\tilde n_\ell}$ parametrized (in most cases linearly) by $\be_\ell$,
	\item a nonlinearity independent of $\be_\ell$, given by $\rho_\ell : \RR^{\tilde n_{\ell}} \rightarrow \RR^{n_{\ell+1}}$
\end{rs}
which we write as
\eql{\label{eq-struc-layer}
	\foralls x_\ell \in \RR^{n_\ell}, \quad
	f_{\ell}(x_\ell,\be_\ell) = \rho_\ell(  B_\ell x_\ell + b_\ell ) \in \RR^{n_{\ell+1}}. 
}
In a fully connected layer, every entry of $B_\ell$ and $b_\ell$ is trainable, so $(B_\ell,b_\ell)=\be_\ell$.
 
A common choice is a pointwise activation $\rho_\ell$, written as $\rho_\ell(z)=(\tilde\rho_\ell(z_k))_k$ with $\tilde\rho_\ell : \RR \rightarrow \RR$ nonlinear. Such an activation preserves dimension, so $n_{\ell+1}=\tilde n_\ell$. Standard examples are the rectified linear unit (ReLU), $\tilde\rho_\ell(s)=\max(s,0)$, and the sigmoid, $\tilde\rho_\ell(s)=\th(s) = (1+e^{-s})^{-1}$.

Interleaving affine maps with nonlinear activations produces increasingly expressive functions $f(\cdot,\be)$.

Training the parameters $\be = (B_\ell,b_\ell)_\ell$ means minimizing~\eqref{eq-erm-param}, for example by stochastic gradient descent. Automatic differentiation computes the gradient at a cost proportional to a forward evaluation, namely $O(\sum_\ell n_\ell\tilde n_\ell)$ for dense layers with pointwise activations. The chain structure gives the backpropagation formula~\eqref{eq-backprop-discr}.

For regression tasks, one usually uses an affine output layer with a squared $\ell^2$ loss $L$. An output nonlinearity is appropriate when the target is constrained, for example to nonnegative values.
 
For classification, a multiclass logistic map~\eqref{eq-multiclass-logitmap} converts the final scores into class probabilities.

For high-dimensional inputs such as images, fully connected layers can require prohibitively many parameters. They also ignore spatial structure and symmetries. Convolutional architectures exploit this structure, as detailed in Section~\ref{sec-cnn}.


\begin{figure}[tbp]
\centering
{\small\sffamily\color{BookMuted}Panel 1\par}\smallskip
\BookDrawing{tikz/deep-learning-fully-connected}
\par\medskip
{\small\sffamily\color{BookMuted}Panel 2\par}\smallskip
\BookDrawing{tikz/deep-learning-convolutional}
\caption{Panel 1: fully connected network.
Panel 2: convolutional neural network.}
\label{fig-deep-fc-cnn}
\end{figure}


                                          
\subsection{Perceptron and Shallow Models}

The logistic classifiers from Sections~\ref{sec-two-class-logit} and~\ref{sec-multiclass-logit} are simple instances of this network formulation.

The binary logistic model~\eqref{eq-two-class-logit-model} is a one-layer ($L=1$) instance of~\eqref{eq-struc-layer}. Omitting the bias, its components are
\eq{
	B_0 x = \dotp{x}{\be}
	\qandq
	\tilde\rho_0(u) = \th(u).
}
The network $f(x,\be) = \th(\dotp{x}{\be})$ can include a bias by appending a constant coordinate to $x$. For labels $y \in \{0,1\}$, it is trained with binary cross-entropy:
\eq{
	L(t,y) = -\log( t^{y} (1-t)^{1-y} ) = -y\log(t)-(1-y)\log(1-t).
}
The resulting empirical-risk objective is convex in the parameters.

For $K$ classes, compute the scores $B_0 x = (\dotp{x}{\be_k})_{k=1}^K$ and apply the normalized exponential map
\eq{
	f(x,\be) = \Nn( ( \exp(\dotp{x}{\be_k}) )_k )
	\qwhereq
	\Nn(u) = \frac{u}{\sum_{k} u_{k}}
} 
For target vectors $y$ on the probability simplex, use the cross-entropy loss
\eq{
	L(t,y) = -\sum_{k=1}^K y_k \log(t_k). 
}

                                          
\subsection{Convolutional Neural Networks}
\label{sec-cnn}

Signals, images, and videos carry spatial structure that can guide the architecture. Convolutional layers respect translations of the spatial grid, subject to boundary conditions, and reuse the same filters at every location.

At depth $\ell$, a convolutional network arranges the activation vector $x_\ell\in\RR^{n_\ell}$ as an array in $\RR^{\bar n_\ell\times d_\ell}$, with $\bar n_\ell$ spatial positions and $d_\ell$ channels, so $n_\ell=\bar n_\ell d_\ell$. The spatial positions typically form a one-, two-, or three-dimensional grid.
 
For an RGB image, $\bar n_0$ is the number of pixels and $d_0=3$.
 
Write $x_\ell = ( (x_\ell)_r[i] )_{r,i}$, where $i \in \{1,\ldots,\bar n_\ell\}$ indexes the spatial position and $r \in \{1,\ldots,d_\ell\}$ the channel.

We require the l